\chapter{Business \& Personal Documents}
\label{ch:business}

This chapter covers five doctypes for professional and personal use:
\texttt{cv}, \texttt{cover-letter}, \texttt{invoice}, \texttt{memo}, and
\texttt{recipe}.  These profiles demonstrate OmniLaTeX's versatility beyond
academic writing.

\section{CV / Resume}
\label{sec:cv}

The \texttt{cv} doctype produces a single-page curriculum vitae with a
compact layout: 10\,pt sans-serif font, empty page style, and minimal
item spacing (\texttt{documentitemspacing=compact}).  It loads
\texttt{fontawesome5}, \texttt{tikz}, and \texttt{adjustbox} for icons
and optional circular profile photos.

\begin{minted}{latex}
\documentclass[doctype=cv]{../../omnilatex}
\title{Jane Doe}
\cvrole{Senior Software Engineer}
\cvcontact{jane@example.com}
\cvlocation{San Francisco, CA}
\cvphone{+1\,555\,012\,3456}
\cvlinks{\href{https://github.com/janedoe}{github.com/janedoe}
    \textbullet{}
    \href{https://janedoe.dev}{janedoe.dev}}
\cvsummary{Experienced engineer with 10+ years in distributed systems.}
\begin{document}
\maketitle
\section{Experience}
\begin{itemize}
    \item \textbf{Staff Engineer} -- Acme Corp, 2020--present
    \item \textbf{Senior Engineer} -- Startup Inc, 2016--2020
\end{itemize}
\section{Education}
\begin{itemize}
    \item \textbf{M.Sc.\ Computer Science} -- University, 2016
\end{itemize}
\end{document}
\end{minted}

\subsection{Custom Commands}

\begin{table}[htbp]
    \centering
    \caption{CV metadata commands}
    \label{tab:cv-commands}
    \begin{tabular}{@{} l p{7.5cm} @{}}
        \toprule
        \textbf{Command} & \textbf{Description} \\
        \midrule
        \texttt{\textbackslash cvrole\{...\}}        & Job title or role \\
        \texttt{\textbackslash cvcontact\{...\}}     & Contact line (email or phone) \\
        \texttt{\textbackslash cvlocation\{...\}}    & City and region \\
        \texttt{\textbackslash cvphone\{...\}}       & Phone number \\
        \texttt{\textbackslash cvlinks\{...\}}       & Free-form links and social URLs \\
        \texttt{\textbackslash cvsummary\{...\}}     & Professional summary (rendered below header rule) \\
        \texttt{\textbackslash cvimage\{...\}}       & Path to a profile photo (rendered as circular crop) \\
        \texttt{\textbackslash cvbulletspacing\{...\}} & Set \texttt{\textbackslash itemsep} for list items \\
        \bottomrule
    \end{tabular}
\end{table}

The \texttt{\textbackslash maketitle} command renders the header with name,
role, location, phone, contact, and links.  When \texttt{\textbackslash cvimage}
is set, a circular photo (2\,cm diameter via TikZ clipping) appears to the
left of the header block.  Each \texttt{\textbackslash section} is followed
by a horizontal rule for visual separation.

\subsection{cv-twopage Variant}

The \texttt{examples/cv-twopage/} example demonstrates a two-page CV using
\texttt{scrbook} instead of \texttt{scrartcl}.  It uses the same
\texttt{cv} doctype commands but allows a longer document with chapter
structure.  See \texttt{examples/cv-twopage/config/document-settings.sty}
for the full configuration.

\section{Cover Letter}
\label{sec:cover-letter}

The \texttt{cover-letter} doctype formats a formal cover letter with 11\,pt
serif font, one-and-a-half line spacing, and a plain page style.  The
\texttt{\textbackslash maketitle} command is disabled; metadata commands are
used to populate the letter's address block manually.

\begin{minted}{latex}
\documentclass[doctype=cover-letter]{../../omnilatex}
\coverletterrecipient{Dr.\ A.\ Smith}
\coverletterposition{Senior Engineer}
\coverlettercompany{Acme Corp}
\coverletteraddress{123 Main St}
\coverlettercity{San Francisco, CA 94102}
\coverletterdate{2025-11-01}
\coverlettersendername{Jane Doe}
\coverlettersenderaddress{456 Oak Ave}
\coverlettersendercity{Berkeley, CA 94701}
\coverlettersenderphone{+1\,555\,012\,3456}
\coverlettersenderemail{jane@example.com}
\begin{document}
% Use the stored metadata to build the letter body
\end{document}
\end{minted}

\begin{table}[htbp]
    \centering
    \caption{Cover letter commands}
    \label{tab:cover-letter-commands}
    \begin{tabular}{@{} l p{7.5cm} @{}}
        \toprule
        \textbf{Command} & \textbf{Description} \\
        \midrule
        \texttt{\textbackslash coverletterrecipient\{...\}}  & Recipient name \\
        \texttt{\textbackslash coverletterposition\{...\}}   & Job position title \\
        \texttt{\textbackslash coverlettercompany\{...\}}    & Company name \\
        \texttt{\textbackslash coverletteraddress\{...\}}    & Recipient street address \\
        \texttt{\textbackslash coverlettercity\{...\}}       & Recipient city and postcode \\
        \texttt{\textbackslash coverletterdate\{...\}}       & Letter date (default: \texttt{\textbackslash today}) \\
        \texttt{\textbackslash coverlettersendername\{...\}} & Sender name \\
        \texttt{\textbackslash coverlettersenderaddress\{...\}} & Sender street address \\
        \texttt{\textbackslash coverlettersendercity\{...\}} & Sender city and postcode \\
        \texttt{\textbackslash coverlettersenderphone\{...\}} & Sender phone \\
        \texttt{\textbackslash coverlettersenderemail\{...\}} & Sender email \\
        \bottomrule
    \end{tabular}
\end{table}

\section{Invoice}
\label{sec:invoice}

The \texttt{invoice} doctype produces a professional invoice with a centered
title, from/to address blocks, and a line-item table.  It uses
\texttt{DIV=10} for generous margins and \texttt{pagestyle=empty}.

\begin{minted}{latex}
\documentclass[doctype=invoice]{../../omnilatex}
\invnumber{INV-2025-0042}
\invdate{2025-10-01}
\invdue{2025-10-31}
\invfrom{Jane Doe Consulting}
\invto{Acme Corporation}
\invfromaddr{456 Oak Ave, Berkeley, CA 94701}
\invtoaddr{123 Main St, San Francisco, CA 94102}
\invcurrency{USD}
\invtaxrate{10}
\invpaymentterms{Net 30}
\begin{document}
\maketitle
\begin{tabular}{@{} p{6cm} r r r @{}}
    \toprule
    \textbf{Description} & \textbf{Qty} & \textbf{Price} & \textbf{Total} \\
    \midrule
    \invoiceitem{Backend Development}{40}{150}
    \invoiceitem{Code Review}{10}{200}
    \bottomrule
\end{tabular}
\end{document}
\end{minted}

\begin{table}[htbp]
    \centering
    \caption{Invoice commands}
    \label{tab:invoice-commands}
    \begin{tabular}{@{} l p{7.5cm} @{}}
        \toprule
        \textbf{Command} & \textbf{Description} \\
        \midrule
        \texttt{\textbackslash invnumber\{...\}}       & Invoice number (default: INV-0001) \\
        \texttt{\textbackslash invdate\{...\}}         & Issue date \\
        \texttt{\textbackslash invdue\{...\}}          & Due date \\
        \texttt{\textbackslash invfrom\{...\}}         & Sender company name \\
        \texttt{\textbackslash invto\{...\}}           & Client name \\
        \texttt{\textbackslash invfromaddr\{...\}}     & Sender address \\
        \texttt{\textbackslash invtoaddr\{...\}}       & Client address \\
        \texttt{\textbackslash invcurrency\{...\}}     & Currency code (default: USD) \\
        \texttt{\textbackslash invtaxrate\{...\}}      & Tax rate percentage (default: 0) \\
        \texttt{\textbackslash invpaymentterms\{...\}} & Payment terms (default: Net 30) \\
        \texttt{\textbackslash invnotes\{...\}}        & Additional notes \\
        \bottomrule
    \end{tabular}
\end{table}

The \texttt{\textbackslash invoiceitem\{desc\}\{qty\}\{price\}} command
emits a table row with description, quantity, unit price, and computed
total, followed by a horizontal rule.

\section{Memo}
\label{sec:memo}

The \texttt{memo} doctype formats a business memorandum with 12\,pt serif
font and single line spacing.  The \texttt{\textbackslash maketitle}
command renders a MEMORANDUM header with ruled lines and labelled fields.

\begin{minted}{latex}
\documentclass[doctype=memo]{../../omnilatex}
\memoto{Engineering Team}
\memofrom{Jane Doe, CTO}
\memosubject{Q4 Development Roadmap}
\memodate{2025-10-15}
\memocc{CFO Office}
\memoref{MEMO-2025-042}
\begin{document}
\maketitle
This memo outlines the development priorities for Q4\ldots
\end{document}
\end{minted}

\begin{table}[htbp]
    \centering
    \caption{Memo commands}
    \label{tab:memo-commands}
    \begin{tabular}{@{} l p{7.5cm} @{}}
        \toprule
        \textbf{Command} & \textbf{Description} \\
        \midrule
        \texttt{\textbackslash memoto\{...\}}      & Recipient \\
        \texttt{\textbackslash memofrom\{...\}}    & Sender \\
        \texttt{\textbackslash memosubject\{...\}} & Subject line \\
        \texttt{\textbackslash memodate\{...\}}    & Date (default: \texttt{\textbackslash today}) \\
        \texttt{\textbackslash memocc\{...\}}      & CC list (omitted when empty) \\
        \texttt{\textbackslash memobcc\{...\}}     & BCC list (omitted when empty) \\
        \texttt{\textbackslash memoref\{...\}}     & Reference number (omitted when empty) \\
        \bottomrule
    \end{tabular}
\end{table}

The header renders as \textbf{TO}, \textbf{FROM}, \textbf{DATE}, and
\textbf{RE} fields with a 1.5\,cm label width.  \texttt{REF} and
\texttt{CC} fields are only displayed when non-empty.

\section{Recipe}
\label{sec:recipe}

The \texttt{recipe} doctype formats cooking recipes with metadata, an
ingredients list, step-by-step instructions, and optional chef's notes.
It uses \texttt{parskip=half} and single line spacing.

\begin{minted}{latex}
\documentclass[doctype=recipe]{../../omnilatex}
\title{Classic Margherita Pizza}
\author{Jane Doe}
\recipeMeta{prep}{30 minutes}
\recipeMeta{cook}{12 minutes}
\recipeMeta{servings}{4}
\recipeMeta{difficulty}{Easy}
\recipeMeta{cuisine}{Italian}
\recipeMeta{category}{Main Course}
\recipeMeta{calories}{280}
\begin{document}
\maketitle
\begin{ingredients}
    \item 500g bread flour
    \item 7g active dry yeast
    \item 200ml warm water
\end{ingredients}
\begin{instructions}
    \item Combine flour, yeast, and salt in a large bowl.
    \item Add warm water and mix until a dough forms.
\end{instructions}
\recipeNote{For a crispier crust, preheat the oven to 250\textdegree C.}
\end{document}
\end{minted}

\subsection{Metadata Commands}

\begin{table}[htbp]
    \centering
    \caption{Recipe metadata commands}
    \label{tab:recipe-commands}
    \begin{tabular}{@{} l p{7.5cm} @{}}
        \toprule
        \textbf{Command} & \textbf{Description} \\
        \midrule
        \texttt{\textbackslash recipeprep\{...\}}       & Preparation time (default: --) \\
        \texttt{\textbackslash recipecook\{...\}}       & Cooking time (default: --) \\
        \texttt{\textbackslash recipeservings\{...\}}   & Number of servings (default: --) \\
        \texttt{\textbackslash recipedifficulty\{...\}} & Difficulty level (default: --) \\
        \texttt{\textbackslash recipecuisine\{...\}}    & Cuisine type (default: --) \\
        \texttt{\textbackslash recipecategory\{...\}}   & Meal category (default: --) \\
        \texttt{\textbackslash recipecalories\{...\}}   & Calories per serving (default: --) \\
        \bottomrule
    \end{tabular}
\end{table}

All metadata fields can also be set via the helper
\texttt{\textbackslash recipeMeta\{key\}\{value\}}, which expands to
\texttt{\textbackslash recipe\textit{key}\{value\}}.  The metadata is
rendered in a tcolorbox with a blue frame below the title.

\subsection{Environments and Helpers}

\begin{itemize}
    \item \texttt{ingredients} -- A bulleted list with compact spacing
          (\texttt{leftmargin=1.5em}, \texttt{itemsep=2pt}).
    \item \texttt{instructions} -- A numbered list with slightly wider
          spacing (\texttt{leftmargin=2em}, \texttt{itemsep=4pt}).
    \item \texttt{\textbackslash recipeNote\{...\}} -- A tcolorbox with an
          orange frame, titled ``Chef's Note'', for tips and variations.
\end{itemize}
